\appendix
\section{Phụ lục --- Khai báo phạm vi dùng AI tạo sinh}
\label{sec:phu-luc-ai}

Đề yêu cầu khai báo phạm vi dùng AI tạo sinh trong phụ lục; từng thành viên có thể
bị vấn đáp; không dùng bản AI miễn phí \cite[tr.~2, mục II]{de}. Phụ lục này khai
báo trung thực, kể cả những phần AI làm phần lớn công việc.

Nguồn chi tiết theo từng phiên là \texttt{docs/progress/ai-use.md}, được cập nhật
liên tục trong suốt dự án.

\subsection{Công cụ và trạng thái xác minh}

Công cụ: một trợ lý lập trình AI chạy trong workspace của nhóm, có hỗ trợ soạn tài
liệu.

\textbf{Điểm cần nêu rõ:} nhóm \textbf{chưa xác minh} loại tài khoản hay gói dịch
vụ AI đang dùng có phù hợp quy định ``không dùng bản AI miễn phí'' hay không. Nhóm
không diễn giải quy định này thành cấm dùng Colab hay Kaggle bản miễn phí, vì đó là
nền tảng tính toán, không phải công cụ AI tạo sinh. Việc xác minh tài khoản là hạng
mục còn mở.

\subsection{Phạm vi AI đã làm}

\subsubsection{Tài liệu, kế hoạch và hồ sơ}

AI kiểm kê Git và tài liệu; đọc và đối chiếu các PDF trong \texttt{references/};
soạn ma trận yêu cầu, đề cương, kế hoạch dữ liệu, hướng dẫn kho mã, bối cảnh bàn
giao, kế hoạch phối hợp nhóm và lịch chuẩn bị vấn đáp; tạo mẫu LaTeX, bộ lọc Pandoc
và script sinh/biên dịch PDF; tổ chức lại README.

\subsubsection{Mã Phần A và bộ kiểm}

\begin{itemize}[leftmargin=2em]
\item \textbf{A.1:} AI viết \texttt{src/attention.py}, notebook dẫn xuất và bộ
kiểm; thư viện chuẩn của PyTorch chỉ dùng ở phía đối chiếu. Chạy trên CPU cục bộ.
\item \textbf{A.2:} AI viết LM dùng decoder A.1, mã đọc dữ liệu có kiểm quyền, vòng
huấn luyện, tính PPL, lưu/khôi phục checkpoint, notebook và bộ kiểm. AI không chạy
Kaggle thay nhóm.
\item \textbf{A.3:} AI triển khai lớp LoRA, notebook và bộ kiểm CPU; chạy đối chiếu
số tham số và phép hợp nhất sau cập nhật khác $0$.
\item \textbf{A.4:} AI triển khai hàm DPO, notebook và bộ kiểm CPU bằng số minh
họa; đối chiếu mất mát và gradient.
\end{itemize}

AI cũng đã phát hiện, tái hiện và sửa một lỗi trong script corpus (theo redirect
trước khi kiểm URL đích), kèm kiểm thử hồi quy.

\subsubsection{Phần B và Phần C}

\begin{itemize}[leftmargin=2em]
\item \textbf{Phần B:} AI đối chiếu sáu yêu cầu với đề tr.~18, đọc các phần liên
quan của \cite{attention}, \cite{lora}, \cite{dpo}, và \textbf{soạn toàn bộ bản
nháp LaTeX B.1--B.6}: phát biểu, giả thiết, các bước chứng minh, diễn giải và ví dụ
minh họa. Ở một lượt rà sau đó, AI tìm ra lỗ hổng số chiều lẻ của B.3 và soạn phần
bổ sung ở Mục~\ref{sec:b3-odd}.
\item \textbf{C.1:} theo yêu cầu của nhóm, AI sinh ma trận Q/K/V với seed
\texttt{577893521}, lưu đầu vào tái lập, soạn phép tính tay từng token và đối chiếu
với mã A.1. Việc giảng viên không cung cấp ma trận là thông tin do nhóm báo.
\item \textbf{C.2:} AI dẫn xuất số tham số và MAC theo kích thước tensor của
\texttt{DecoderBlock}, rồi chạy notebook CPU đối chiếu \texttt{numel()} và FLOP do
profiler ước lượng.
\item \textbf{C.3:} AI đối chiếu cấu hình Qwen khóa revision, mã Qwen2 trong
Transformers và giả định double quantization của \cite{qlora}; viết notebook và
chạy trên CPU để kiểm số tham số bằng lớp A.3.
\item \textbf{C.4:} AI tự đặt bốn log-xác suất minh họa, trình bày phép thế số, và
kiểm mất mát/gradient với A.4 và autograd.
\end{itemize}

\subsubsection{Phần D}

\begin{itemize}[leftmargin=2em]
\item \textbf{D.1:} AI truy tìm văn bản công khai, đối chiếu số hiệu, ngày, phạm
vi, phụ lục và căn cứ quyền tác giả; lưu URL, thời điểm và SHA-256; kiểm chữ ký CMS
về toàn vẹn nội dung trên bốn PDF Hà Nội. AI không đăng nhập, không lấy hồ sơ công
dân, không tự phê duyệt nguồn.
\item \textbf{D.2:} AI soạn recipe từ chính văn, kiểm SHA-256 và đoạn trích, triển
khai CLI chọn mẫu phân tầng và cổng xuất SFT, cài công cụ kiểm tách
\texttt{USER\_TEST}. AI biên soạn 20 QA nghiên cứu và đối chiếu 21 đoạn trích.
\item \textbf{D.3:} AI chuẩn bị bộ nạp release D.2, mã QLoRA, notebook Kaggle,
script đóng gói và kế hoạch so cấu hình trên validation; chạy kiểm phần mềm trên
CPU.
\item AI soạn bộ tài liệu đồng ý, khảo sát, quy trình riêng tư và rubric/phiếu
chấm.
\end{itemize}

\subsubsection{Báo cáo này}

AI soạn cấu trúc và toàn bộ bản nháp của báo cáo này từ các hiện vật đã có trong
kho mã, kèm datasheet và model card. AI đối chiếu lại từng con số với notebook, tệp
tài liệu nguồn và ma trận yêu cầu; chạy bộ kiểm tự động và ghi kết quả thực tế. AI
không tạo kết quả thực nghiệm mới và không đổi trạng thái của hạng mục nào.

\subsection{Phạm vi AI \emph{không} làm}

Danh sách này quan trọng ngang danh sách trên:

\begin{enumerate}[leftmargin=2em]
\item \textbf{Không thực hiện khảo sát con người} và không tạo bản ghi khảo sát.
Chưa mời hay thu dữ liệu của bất kỳ người tham gia nào.
\item \textbf{Không gán nhãn ưu tiên thật} và không tạo mức đồng thuận.
\item \textbf{Không điền phiếu kiểm thủ công.} Trong vòng rà pilot, AI chạy
validator ngoại tuyến và lưu báo cáo riêng, nhưng số QA được người duyệt vẫn giữ
bằng $0$. Nhận xét của AI \emph{không} được tính vào quota kiểm thủ công $\ge20\%$.
\item \textbf{Không chạy huấn luyện thay nhóm} trên Colab hay Kaggle, không tạo
dataset/notebook trên tài khoản nền tảng của nhóm.
\item \textbf{Không tự phê duyệt quyền dùng nguồn} và không ký xác nhận thay người
có thẩm quyền. Kết quả audit kỹ thuật không được dùng làm attestation về quyền.
\item \textbf{Không thay bộ Q/K/V chính thức của Phần C} --- bộ minh họa luôn được
ghi nhãn là minh họa.
\item \textbf{Không tạo kết quả đánh giá} hay số đo GPU.
\end{enumerate}

\subsection{Dữ liệu đầu vào cho AI}

Yêu cầu của nhóm; sáu PDF có sẵn trong workspace (PDF đề gốc và năm bài báo); các
trang web công khai đã ghi URL và thời điểm trong nhật ký nguồn; và những thông số
phiên Kaggle/Colab \emph{do thành viên nhóm cung cấp}. Các thông số do thành viên
cung cấp được ghi là thông tin người dùng cung cấp, \emph{không} nhận là kết quả do
AI đo.

\subsection{Trách nhiệm của nhóm}

Những việc sau thuộc trách nhiệm thành viên và \emph{không} được chuyển cho AI:

\begin{itemize}[leftmargin=2em]
\item Đọc lại nguồn và xác minh các giả thiết.
\item Lựa chọn phạm vi dự án và mô hình.
\item \textbf{Hiểu và giải thích được toàn bộ mã và chứng minh khi vấn đáp.} Mỗi
thành viên tự giải lại các chứng minh của Phần~B và chạy lại các phép kiểm, không
dựa vào việc AI đã viết chúng.
\item Thực hiện kiểm định thủ công, khảo sát, gán nhãn ưu tiên và đánh giá mù.
\item Ghi nhận sửa đổi do thành viên thực hiện, công cụ/phiên bản và phạm vi AI ở
các phiên sau.
\end{itemize}

\textbf{Trạng thái tại thời điểm viết:} nhóm chưa ghi tên người kiểm duyệt nội dung
cuối vì chưa có xác nhận. Phần~B vẫn mang nhãn bản nháp chưa được nhóm duyệt
(Mục~\ref{sec:phan-b}). Văn bản do AI soạn \textbf{không} được coi là nguồn học
thuật thay cho năm bài báo và đề gốc.
